% ============================================================
%  第六章  辛普森悖论、混杂与可压缩性（第一部分：6.1—6.3）
% ============================================================
\chapter{辛普森悖论、混杂与可压缩性}

\begin{flushright}
\kaishu\small
直面悖论者，即把自己暴露于现实。\\[2pt]
——迪伦马特（Friedrick Durrenmatt，1962）
\end{flushright}

\section*{序言}
\addcontentsline{toc}{section}{序言}

混杂是从经验数据中阐明因果推断的最根本障碍之一。因此，在严重依赖因果推断的领域——流行病学、计量经济学、生物统计学与社会科学——已写就或谈论的许多内容，都以混杂的考量为基础。然而，除随机实验的标准分析之外，多数统计学教科书对这一主题几乎不置一词。原因很简单：混杂是一个因果概念，因而无法在标准统计模型中表达。一旦尝试作形式化的统计分析，往往就导致混乱或繁复，使这一主题对非专家而言极难理解，更不用说掌握。

我写这本书的主要目标之一，就是看到这些混乱得到澄清——看到涉及混杂控制的问题被化归为简单的数学程式。第三章引入的数学技术，确实结晶为若干简单的图形程式：用以检测混杂的存在，并识别为获得无混杂效应估计而应当控制的变量。本章则处理我们试图用统计准则来定义并控制混杂时所遇到的困难。

我们从剖析辛普森悖论的有趣历史入手（\S 6.1），并以它作为放大镜，考察一代代统计学家在试图用统计学语言捕捉因果概念时所遇到的困难。在 \S 6.2 与 \S 6.3，我们考察用仅基于频率数据与可测统计关联的统计准则来取代混杂之因果定义的可行性。我们将表明，尽管这种取代一般不可行（\S 6.3），但某一类无混杂条件——称为“稳定的”——仍可给出统计的或半统计的刻画（\S 6.4）。这一刻画导致若干操作检验，它们类似于可压缩性检验，可以提醒研究者注意某给定效应估计中存在的不稳定或偏倚（\S 6.4.3）。最后，\S 6.5 澄清可压缩性与无混杂之别、混杂因子与混杂之别，以及表示混杂问题的结构进路与可交换性进路之别。

\section{辛普森悖论剖析}

本书已两次简要讨论过被称为辛普森悖论的逆转效应：第一次是在协变量选择问题（\S 3.3）的语境中，第二次是在直接效应定义（\S 4.5.3）的语境中。本节分析逆转效应为何曾（且仍）被视为悖论，以及它的解决为何来得如此之迟。

\subsection{一个非悖论的故事}

辛普森悖论（Simpson 1951; Blyth 1972）最早由 Pearson 于 1899 年遇到（Aldrich 1995），指的是这样一种现象：事件 $C$ 在给定总体 $p$ 中增大 $E$ 的概率，同时又减小 $E$ 在 $p$ 的每个子总体中的概率。换言之，若 $F$ 与 $\neg F$ 是描述两个子总体的两种互补属性，我们完全可能遇到下述不等式
\begin{gather}
P(E\mid C) > P(E\mid \neg C), \tag{6.1}\\
P(E\mid C,F) < P(E\mid \neg C,F), \tag{6.2}\\
P(E\mid C,\neg F) < P(E\mid \neg C,\neg F). \tag{6.3}
\end{gather}

尽管这种次序逆转或许不会让概率论的学生感到惊讶，但一旦赋予其因果解释，它便成为悖论。例如，若把 $C$（蕴含“原因”）与服用某种药物相联系，把 $E$（蕴含“效应”）与康复相联系，把 $F$ 与“女性”相联系，那么在 \textup{(6.2)}—\textup{(6.3)} 的因果解释下，该药物似乎对男性与女性皆有害，却对整个总体有益（式 \textup{(6.1)}）。直觉判定这样的结果不可能，而这一判定是正确的。

图 6.1 的三张表以数值方式呈现了辛普森的逆转。可以看到，总体而言，服用药物（$C$）者的康复率为 50\%，高于对照组的 40\%，因而药物治疗显然可取。然而，当我们分别查看男性与女性的表时，未治疗者的康复率在男性与女性中都比治疗者高 10 个百分点。

\begin{figure}[htbp]
\centering
\begin{tabular}{@{}p{3.35cm} r r r r@{}}
\toprule
\textbf{（a）合并} & $E$ & $\neg E$ & 合计 & 康复率\\
\midrule
服药（$C$） & 20 & 20 & 40 & 50\%\\
不服药（$\neg C$） & 16 & 24 & 40 & 40\%\\
\midrule
合计 & 36 & 44 & 80 & \\
\bottomrule
\end{tabular}

\par\vspace{2.0em}
\begin{tabular}{@{}p{3.35cm} r r r r@{}}
\toprule
\textbf{（b）男性} & $E$ & $\neg E$ & 合计 & 康复率\\
\midrule
服药（$C$） & 18 & 12 & 30 & 60\%\\
不服药（$\neg C$） & 7 & 3 & 10 & 70\%\\
\midrule
合计 & 25 & 15 & 40 & \\
\bottomrule
\end{tabular}

\par\vspace{2.0em}
\begin{tabular}{@{}p{3.35cm} r r r r@{}}
\toprule
\textbf{（c）女性} & $E$ & $\neg E$ & 合计 & 康复率\\
\midrule
服药（$C$） & 2 & 8 & 10 & 20\%\\
不服药（$\neg C$） & 9 & 21 & 30 & 30\%\\
\midrule
合计 & 11 & 29 & 40 & \\
\bottomrule
\end{tabular}
\caption{服药（$C$）与对照（$\neg C$）在男性、女性以及合并样本中的康复率。}
\label{fig:6.1}
\end{figure}

对本书读者而言，辛普森悖论的解释应当已经清楚，因为我们曾不遗余力地区分“看”与“做”。概率演算中的取条件算子代表证据性条件句“已知我们看到”，而 $\dox{\cdot}$ 算子则用以表示因果条件句“已知我们去做”。据此，不等式
\[
P(E\mid C) > P(E\mid \neg C)
\]
并非关于 $C$ 对 $E$ 有正向效应的陈述——那应正确写作
\[
P(E\mid \dox{C}) > P(E\mid \dox{\neg C}),
\]
而是关于 $C$ 是 $E$ 之正证据的陈述；这种证据可能来自同时导致 $C$ 与 $E$ 的虚假混杂因子。在我们的例子里，药物之所以在总体上显得有益，是因为男性（无论是否服药都）比女性更常康复，同时也比女性更可能服用该药。确实，若发现一位性别未知的服药病人（$C$），我们大可以推断此人更可能是男性，因而更可能康复——这与 \textup{(6.1)}—\textup{(6.3)} 完全吻合。

处理这类潜在混杂因子的标准方法是把它们“固定”\footnote{“把 $F$ 固定”与“对 $F$ 作控制”这两种说法，哲学家（如 Eells 1991）与统计学家（如 Pratt and Schlaifer 1988）都在使用，但它们蕴含外部干预，因而可能引起误解。在统计分析中，我们所能做的只是通过考察 $F$ 取相同值的个案来“模拟”把 $F$ 固定——也就是对 $F$ 与 $\neg F$ 取条件——这一操作我称为“对 $F$ 作调整”。}，即对任何可能同时导致 $C$ 与 $E$ 的因素取条件。在我们的例子里，若认为身为男性（$\neg F$）既是康复（$E$）的原因、又是服药（$C$）的原因，则药物的效应须对男性与女性分别评估（如 \textup{(6.2)}—\textup{(6.3)}），再据此取平均。于是，假定 $F$ 是唯一的混杂因子，\textup{(6.2)}—\textup{(6.3)} 恰当地表示了药物在相应总体中的效力，而 \textup{(6.1)} 仅仅表示在缺乏性别信息时它的证据权重，悖论便消解了。

\subsection{一段统计学的痛苦史}

至此，我们描述的是现代因果性研究者所理解（或应当理解）的悖论（例如见 Cartwright 1983;\footnote{Cartwright 却主张，第三因素 $F$ 应当“被固定”，当且仅当 $F$ 与 $E$ 有因果关联（p.~37）；正确的（后门）准则要更复杂一些（见定义 3.3.1）。} Holland and Rubin 1983; Greenland and Robins 1986; Pearl 1993b; Spirtes et al. 1993; Meek and Glymour 1994）。然而，多数统计学家不愿接受“辛普森悖论源自因果考量”这一想法。普遍的态度是：逆转是真实的、令人不安的，因为它确实出现在数字中，且确实可能把统计学家引向错误结论。若某事物是真实的，它就不可能是因果的，因为因果性是一种定义不清的心智构造。于是，该悖论必定是一种统计现象，可以用统计分析的工具检测、理解并回避。《统计学科学百科全书》就严厉警告我们辛普森悖论所潜伏的危险，却只字不提“原因”或“因果性”（Agresti 1983）。《生物统计学百科全书》（Dong 1998）与《医学科学统计学剑桥辞典》（Everitt 1995）也持同样的观念。

据我所知，统计文献中只有两篇文章明确把辛普森逆转的奇特之处归因于因果解释。第一篇是 Pearson 等 (1899)，该现象的发现\footnote{Pearson 等 (1899) 与 Yule (1903) 报告的是悖论的一个较弱版本，其中 \textup{(6.2)}—\textup{(6.3)} 以等号成立。逆转本身后来由 Cohen and Nagel (1934, p.~449) 发现。}在该文中是这样表述的：

\begin{quote}
\kaishu
对那些执意把一切相关都看作因果的人而言，只要取两个亲缘相近的种族作人为混合，就能在两个毫不相关的性状 $A$ 与 $B$ 之间制造出相关——这一事实必定令他们大为震惊。（p.~278）
\end{quote}

在 Pearson 毕生运动的影响下，统计学家只要可能就避谈因果；半个多世纪里，逆转现象被当作 $2\times 2$ 表的一种奇特数学性质来处理，其因果起源被剥离殆尽。

最后，Lindley and Novick (1981) 从新的角度分析了这一问题，并建立了与因果性的第二次公开联系：

\begin{quote}
\kaishu
上一段引入了“原因”这一概念。一种可能的做法是使用因果的语言，而不是可交换性或总体识别的语言。我们未选择这样做，也未讨论因果性，因为这一概念虽然被广泛使用，却似乎并未得到良好的定义。（p.~51）
\end{quote}

关于辛普森逆转的历史，令人惊异之处在于：从 Pearson 等直到 Lindley and Novick，众多就该主题写作的作者中，没有一人敢于追问：这一现象为何值得我们注意？它为何引起惊讶？毕竟，取条件后看到概率大小改变是司空见惯的，而看到这种改变转为符号逆转（对这些概率作差与混合）也并不罕见。那么，若非某种被误导却顽固的幻觉在作祟，不等式改变方向又有什么可震惊的呢？

Pearson 明白，震惊源自被扭曲的因果解释；他着手通过统计相关与列联表的棱镜来纠正这一点（见第十章之后的“后记”）。他的门徒把他看得相当认真，有人甚至断言因果不过是相关的一个种（Niles 1922）。在如此否认一切因果直觉之后，研究者往往别无选择，只能把辛普森逆转归咎于数据的某种邪恶特征——一种严谨的研究者理应回避的特征。自 1950 年代以来，关于辛普森逆转之统计侧面的论文已有数十篇；有些讨论效应的大小（Blyth 1972; Zidek 1984），有些确立其消失的条件（Bishop et al. 1975; Whittemore 1978; Good and Mittal 1987; Wermuth 1987），有些甚至提出激烈如用 $P(C\mid E)$ 取代 $P(E\mid C)$ 作为疗效度量的补救方案（Barigelli and Scozzafava 1984）——无论如何，逆转必须被回避。

对该主题的一种典型处理，见于 Bishop、Fienberg 与 Holland (1975) 那本有影响的著作。Bishop 等 (1975, pp.~41--2) 给出一个例子：产前保健量与婴儿存活之间的表观关联，在按参与研究的各诊所分别考察数据时消失了。他们总结道：“如果我们只看这张（合并的）表，就会错误地断定存活与所接受的保健量\emph{相关}（斜体为我所加）。”讽刺的是，在所研究的样本中，存活确实与所接受的保健量相关。Bishop 等想说的是：不加批判地看合并表，我们会错误地断定存活与所接受的保健量有\emph{因果}关系。然而，由于 1970 年代必须回避因果词汇，像 Bishop 等这样的研究者被迫使用“相关”“伴随”之类的统计替代词，于是自然成为这种语言之局限的牺牲品；统计替代词无法表达研究者本想传达的因果关系。

辛普森悖论帮助我们体会这一代统计学家的痛苦与成就。他们受健全因果直觉的驱使，却在文化上被禁止承认它，在数学上又无力表达它；尽管如此，他们仍从枯燥的表中提炼出意义，并使统计方法成为经验科学的标准。但辛普森悖论的辛辣之处，终究是非统计的。

\subsection{因果性对可交换性}

Lindley and Novick (1981) 最先证明了辛普森悖论的非统计性质——不存在任何统计准则，可以警告研究者不要得出错误结论，或指出哪张表代表正确答案。

在贝叶斯决策论的传统中，他们先把注意力转向该现象的实际一面，并大胆追问：一位新病人来了，我们用不用这药？等价地：我们查哪张表，合并表还是分性别表？“表面上的答案是，”Novick (1983, p.~45) 坦承，“当我们知道病人的性别为男，或知道为女时，我们不用治疗；但若性别未知，我们就该用治疗！显然这个结论荒谬。”Lindley and Novick 随后进行了冗长的非形式讨论，结论（与我们 \S 6.1.1 的结论相同）是我们应当查分性别的表，不用该药。

下一步是追问：是否有某些额外的统计信息，一般地能指引我们找到正确的表？Lindley and Novick 对此给出否定回答：他们表明，面对完全相同的数据，我们有时应当作出相反的决定，去查合并表。他们问：假定数字不变，只改变数据背后的故事，设想 $F$ 代表某种受 $C$ 影响的属性——比如低血压，如图 6.2(b) 所示。\footnote{Lindley and Novick (1981) 所用的例子取自农业，未提及 $C$ 与 $F$ 之间的因果关系，但结构与图 6.2(b) 相同。}检视图 6.2(b) 的图，读者应当立刻得出结论：合并表代表我们想要的答案；我们不应对 $F$ 取条件，因为它恰恰位于我们想要评估的那条因果通路上。（等价地说，通过比较处理后血压相同的病人，我们把药物借以促成康复的两条通路之一的效应掩盖了。）

\begin{figure}[htbp]
\centering
\begin{tikzpicture}[causal,scale=0.92,every node/.style={transform shape}]
% ---------- (a) ----------
\node[vtx] (a1) at (0.5,2.1) {$C$};
\node[vtx] (a2) at (2.9,2.1) {$F$};
\node[vtx] (a3) at (1.5,0.2) {$E$};
\draw[ed] (a2)--(a1);
\draw[ed] (a1)--(a3);
\draw[ed] (a2)--(a3);
\node[font=\footnotesize,above=0.5mm of a1] {处理};
\node[font=\footnotesize,right=1.2mm of a2] {性别};
\node[font=\footnotesize,left=1.2mm of a3] {康复};
\node[font=\small] at (1.7,-0.9) {(a)};
% ---------- (b) ----------
\node[vtx] (b1) at (5.4,2.1) {$C$};
\node[vtx] (b2) at (7.8,2.1) {$F$};
\node[vtx] (b3) at (6.4,0.2) {$E$};
\draw[ed] (b1)--(b2);
\draw[ed] (b1)--(b3);
\draw[ed] (b2)--(b3);
\node[font=\footnotesize,above=0.5mm of b1] {处理};
\node[font=\footnotesize,right=1.2mm of b2,align=left] {血压};
\node[font=\footnotesize,left=1.2mm of b3] {康复};
\node[font=\small] at (6.6,-0.9) {(b)};
% ---------- (c) ----------
\node[vtx] (c1) at (10.2,2.1) {$C$};
\node[vtx] (c2) at (12.2,1.8) {$F$};
\node[vtx] (c3) at (10.9,0.2) {$E$};
\node[lat,minimum size=2.6mm,inner sep=0pt] (c4) at (12.85,3.05) {};
\node[lat,minimum size=2.6mm,inner sep=0pt] (c5) at (13.25,0.5) {};
\draw[ed] (c1)--(c3);
\draw[edd] (c4) to[bend left=8] (c1);
\draw[edd] (c4) to[bend right=14] (c2);
\draw[edd] (c5) to[bend left=14] (c2);
\draw[edd] (c5) to[bend right=4] (c3);
\node[font=\footnotesize,above=0.5mm of c1] {处理};
\node[font=\footnotesize,left=1.2mm of c3] {康复};
\node[font=\small] at (11.8,-0.9) {(c)};
\end{tikzpicture}
\caption{三个能够生成图 6.1 之数据的因果模型。模型 (a) 要求使用分性别的表，而 (b) 与 (c) 要求使用合并表。}
\label{fig:6.2}
\end{figure}

当两个因果模型生成相同的统计数据（图 6.2(a) 与 (b) 是可观测等价的），而在其中一个模型中我们决定用药、在另一个中决定不用药，那么显然，我们的决定由因果考量而非统计考量驱动。有些读者可能猜测时间信息参与了这个决定，指出性别在治疗之前确定、血压在治疗之后确定。但事实并非如此；图 6.2(c) 表明，$F$ 可以发生在 $C$ 之前或之后，而正确的决定仍应是查合并表（即不对 $F$ 取条件，这从后门准则可以看出）。

我们刚刚以例子演示了 \S 6.1.1 已知的道理——即每一个与行动效应有关的问题都必须由因果考量来决定；仅凭统计信息是不够的。此外，选择据以作出决定的正确表格，是协变量选择问题的一个特例，而该问题已在 \S 3.3 用因果演算给出一般解。另一方面，Lindley and Novick 止步于这一认识之前，而把两例之差归因于一个称为“可交换性”的元统计\footnote{所谓“元统计的”，我指一种本身不能从统计数据中识别、却用以评判某种统计方法是否恰当的准则。}概念，它由 DeFinetti (1974) 首先提出。

可交换性关乎这样一个问题：为对某个个体单元作出预测，应选择哪个合适的参考类或子总体。例如，保险公司希望用与某新客户特征最相近的一类人的死亡记录，来估计该客户的平均余命。De Finetti 给这个问题一个形式化的转折，把关于相似性的判断翻译为关于概率的判断。按照这一准则，一个（第 $n+1$ 个）单元在属性 $X$ 上相对于另外 $n$ 个单元构成的组是可交换的，若联合概率分布 $P(X_1,\dots,X_n,X_{n+1})$ 在置换下不变。对 De Finetti 而言，这种不变性如何建立，是一个次要的心理学问题；要点在于把这一心理活动的目标表述为数学表达式，使之能以科学术语交流与讨论。正是这一概念，Lindley and Novick 试图引入辛普森逆转现象，并希望借此表明：在 $F=$ 性别的例子里，合适的子总体是男性与女性；而在 $F=$ 血压的例子里，应考虑的是病人的整个总体。

Lindley and Novick 文章的读者很快就会觉察到，尽管这两位作者用可交换性与子总体的话语来装点讨论，他们实际所做的，是为其直觉结论提供非形式的因果论证。Meek and Glymour (1994) 敏锐地指出，Lindley and Novick 关于可交换性的讨论中唯一可理解的部分，正是基于因果考量的部分；这表明“要理解何时应当、何时不应当假定可交换性，必须对因果信念与概率的相互作用作出明确的说明”（Meek and Glymour 1994, p.~1013）。

实情确实如此；实验研究中的可交换性取决于对生成数据之机制的因果理解。判定一个新单元的响应是否应以某组单元此前的响应为准，取决于这样一个问题：我们打算让新单元经受的实验条件，是否等同于观察该组时所处的条件。我们不能用合并表（图 6.1(a)）来确定一位（性别未知的）新病人的响应，原因在于实验条件已经改变：该组是在病人自行选择治疗的条件下被研究的，而新病人将被命令给予治疗，也许是违背其自然倾向的。因此，新实验将改变某一机制；若不首先就有关概率是否会对这种改变保持不变作出因果假设，就不可能作出任何可交换性判断。我们之所以能在图 6.2(b) 的血压例子中使用合并表，是因为在该设定下，被改变的治疗选择机制被假定对条件概率 $P(E\mid C)$ 没有影响；也就是说，$C$ 被假定为外生的。（这从图中没有任何后门路径可以清楚地看出。）

请注意，若下一位病人就是所研究组中的一员（假设治疗与效应可以复制，且下一位病人的性别与身份未知），同样的考量仍然成立；若我们把样本置于新的实验条件下，从总体中随机抽取的样本与该总体就不是“可交换的”。必须考虑因果机制的改变，才能判定可交换性在新环境下是否成立。而一旦考虑了因果机制，就不再需要单独作可交换性判断了。

但 Lindley and Novick 本可以直接谈论因果关系来表达想法，为何却选择如此曲折地（经由可交换性）言说？他们对这个问题作了部分回答：“（因果性）虽然被广泛使用，却似乎并未得到良好的定义。”人们自然会问：可交换性怎么可能比据以判断它的那些考量更“定义良好”呢？只有当我们考虑 1981 年统计学家可用的数学工具时，才能理解这个答案。当 Lindley and Novick 写道因果性定义不清时，他们真正的意思是：因果性无法用他们所习惯的任何数学形式写下来。路径图、结构方程与 Neyman–Rubin 记号作为数学语言的潜力，在 1981 年普遍未被认识，其原因见 \S 5.1 与 \S 7.4.3。的确，假如 Lindley and Novick 想用因果术语传达想法，他们甚至无法用数学表达“性别不受药物影响”这一简单而关键的事实，更不用说从该事实推出不那么明显的真理了。\footnote{Lindley and Novick (1981, p.~50) 确实尝试过用概率记号表达这一事实。但由于手头没有 $\dox{\cdot}$ 算子，他们不恰当地写了 $P(F\mid C)$ 而非 $P(F\mid \dox{C})$，并无力地论证我们应把 $P(F\mid C)$ 与 $P(F)$ 等同：“你也可以判断，使用治疗或对照的决定不受未知性别的影响，因而 $F$ 与 $C$ 独立。”奇怪的是，这一决定同样不受未知血压的影响；然而，若我们在图 6.2(b) 的例子中写 $P(F\mid C)=P(F)$，就会得到错误的结果。}他们唯一熟悉的正式语言是概率演算，但正如我们已多次见到的，若不加以恰当扩展，这种演算无法充分处理因果关系。

所幸，过去十年发展出的数学工具，使我们能够以更系统、更友好的方式解决辛普森悖论。

\subsection{一个被解决的悖论（或：人是何种机器？）}

悖论如同视错觉，常被心理学家用来揭示心智的内部运作，因为悖论源自（并放大）隐含假设集之间的潜伏冲突。就辛普森悖论而言，我们面对的是 (i) 因果关系受概率演算定律支配这一假设，与 (ii) 驱动我们因果直觉的那组隐含假设之间的冲突。第一个假设告诉我们 \textup{(6.1)}—\textup{(6.3)} 中的三个不等式是相容的，甚至为我们提供了一个概率模型来证实这一断言（图 6.1）。第二个假设告诉我们：不可能存在一种神奇的药物，既对男性与女性都有害，又同时对整个总体有益。

要解决这一悖论，我们必须或者 (a) 表明我们的因果直觉是误导的或不融贯的，或者 (b) 否认“因果关系受标准概率演算定律支配”这一前提。读者此时必定已经猜到，我们将选择第二个方案；我们在此处以及本书其余部分的立场是：因果性受它自己的逻辑支配，而这种逻辑需要对概率演算作重大扩展。我们仍有责任阐明支配因果直觉的逻辑，并形式地表明：这种逻辑排除了上述神奇药物的存在。

$\dox{\cdot}$ 算子的逻辑恰好适用于此。让我们先把“神奇药物 $C$ 对男性与女性都有害”这一陈述翻译为因果演算中的形式陈述：
\begin{gather}
P(E\mid \dox{C},F) < P(E\mid \dox{\neg C},F), \tag{6.4}\\
P(E\mid \dox{C},\neg F) < P(E\mid \dox{\neg C},\neg F). \tag{6.5}
\end{gather}
我们需要证明 $C$ 必定对整个总体有害；也就是说，不等式
\begin{equation}
P(E\mid \dox{C}) > P(E\mid \dox{\neg C}) \tag{6.6}
\end{equation}
必须被证明与我们所知的药物与性别之关系不相容。

\begin{theorem}[确定事件原则]\footnote{Savage (1954, p.~21) 把确定事件原则提出来作为（关于行动的）偏好的基本公设，暗中假定了定理中的“不变”条款。Blyth (1972) 利用这一疏漏构造了一个表面上的反例。定理 6.1.1 表明，确定事件原则不必作为独立公设陈述——它从“行动是结构方程（或机制）之修改者”这一语义中逻辑地推出。反事实分析见 Gibbard and Harper (1976)。注意，“不变”条款是概率性的；只要各子总体的相对规模保持不变，它允许行动改变个体单元的归类。}
若一个行动 $C$ 在每一个子总体中都增大事件 $E$ 的概率，则它也必定增大 $E$ 在整个总体中的概率，前提是该行动不改变各子总体的分布。
\end{theorem}

\noindent\textbf{定理 6.1.1 的证明}\quad 我们将在本例的语境中证明定理 6.1.1，其中总体被划分为男性与女性；推广到多个划分是直接的。在此语境中，我们需要证明 \textup{(6.4)}—\textup{(6.6)} 中的不等式逆转，与“药物对性别没有影响”这一假设不相容：
\begin{equation}
P(F\mid \dox{C}) = P(F\mid \dox{\neg C}) = P(F). \tag{6.7}
\end{equation}
展开 $P(E\mid \dox{C})$ 并用 \textup{(6.7)}，得到
\begin{equation}
\begin{aligned}
P(E\mid \dox{C}) &= P(E\mid \dox{C},F)\,P(F\mid \dox{C})
   + P(E\mid \dox{C},\neg F)\,P(\neg F\mid \dox{C})\\
&= P(E\mid \dox{C},F)\,P(F) + P(E\mid \dox{C},\neg F)\,P(\neg F).
\end{aligned} \tag{6.8}
\end{equation}
类似地，对 $\dox{\neg C}$ 我们得到
\begin{equation}
\begin{aligned}
P(E\mid \dox{\neg C}) &= P(E\mid \dox{\neg C},F)\,P(F)
   + P(E\mid \dox{\neg C},\neg F)\,P(\neg F).
\end{aligned} \tag{6.9}
\end{equation}
由于 \textup{(6.8)} 右端的每一项都小于 \textup{(6.9)} 中对应的项，我们断定
\[
P(E\mid \dox{C}) < P(E\mid \dox{\neg C}),
\]
从而证明了定理 6.1.1。$\blacksquare$

我们于是看到因果直觉从何而来：我们直觉逻辑中一个显然却关键的假设是，药物不影响性别。这解释了当 $F$ 被解释为受药物影响的中间事件（如图 6.2(b)）时，我们的直觉为何发生如此剧烈的变化。在这种情况下，直觉逻辑告诉我们：找到一种满足 \textup{(6.4)}—\textup{(6.6)} 三个不等式的药物是完全相容的；而且，对 $F$ 作调整将是不恰当的。若 $F$ 受 $C$ 影响，则 \textup{(6.8)} 无法导出，而差值 $P(E\mid \dox{C})-P(E\mid \dox{\neg C})$ 可正可负，取决于 $P(F\mid \dox{C})$ 与 $P(F\mid \dox{\neg C})$ 的相对大小。只要 $C$ 与 $E$ 没有共同原因，我们就应直接从合并表（式 \textup{(6.1)}）而非分 $F$ 的表（式 \textup{(6.2)}—\textup{(6.3)}）评估 $C$ 的效力。

请注意，在我们的分析中，从未假定数据来自随机化研究（即 $P(E\mid \dox{C})=P(E\mid C)$），也从未假定来自平衡研究（即 $P(C\mid F)=P(C\mid \neg F)$）。相反，给定图 6.1 的各表，我们的因果逻辑欣然接受我们面对的是一个不平衡的研究，却仍拒绝接受 \textup{(6.4)}—\textup{(6.6)} 的相容性。同样，人们也能从表中清楚地看到，男性比女性更可能服药；然而，当面对逆转现象时，人们仍会“震惊地”发现，康复率的差异可以通过合并表格而被逆转。

我们可以从这些观察中得出的结论是：人类一般对比率与比例（它们是暂时的）视而不见，而不断寻找因果关系（它们是不变的）。一旦人们把比例解释为因果关系，他们便继续用因果演算、而非比例演算来处理这些关系。假如我们的心智受比例演算支配，图 6.1 根本不会引起任何惊讶，辛普森悖论也绝不会引起它曾引起的关注。



\section{为何不存在混杂的统计检验，为何许多人认为存在，以及为何他们几乎是对的}

\subsection{引论}

混杂是一个简单的概念。若我们着手通过考察两个变量（$X$ 与 $Y$）之间的统计关联来估计其一（$X$）对另一（$Y$）的效应，\footnote{我们将把“效应”“影响”“作用于”这些词限定于其因果解释；而把“关联”一词留给统计依赖关系。}就应当确保该关联不是由所研究的效应之外的因素产生的。虚假关联的出现——例如由于额外变量的影响——称为混杂，因为它易于搅乱我们的解读、并使我们对所研究效应的估计产生偏倚。因此，在概念上我们可以说：当存在第三个变量 $Z$ 同时影响 $X$ 与 $Y$ 时，$X$ 与 $Y$ 是混杂的；这样的变量便称为 $X$ 与 $Y$ 的混杂因子。

这个概念虽然简单，却数十年来一直抗拒形式化处理，而且理由充分：“效应”与“影响”这两个概念本身——相对于它们才必须定义“虚假关联”——一直抗拒数学表述。把效应经验地定义为“在受控随机实验中将会成立的关联”，无法轻易用概率论的标准语言表达，因为该理论处理的是静态条件，不允许我们预测——即使从总体密度函数的完整指定出发——若条件改变（比如从观测研究变为受控研究），将会成立何种关系。这类预测需要以因果假设或反事实假设的形式补充额外信息，而这些信息无法从密度函数中识别出来（见 \S 1.3 与 \S 1.4）。本书所用的 $\dox{\cdot}$ 算子，正是为区分与管理这种额外信息而设计的。

尽管有这些困难，流行病学家、生物统计学家、社会科学家与经济学家\footnote{在计量经济学中，困难集中于“外生性”这一概念（Engle et al. 1983; Leamer 1985; Aldrich 1993），它实质上代表“无混杂”（见 \S 5.4.3）。}仍作过无数尝试，试图用统计术语定义混杂：部分因为统计定义——不涉及“效应”或“影响”这类理论词项——可以用常规数学形式表达；部分因为这类定义可能导向对混杂的实用检验，从而提醒研究者注意可能的偏倚与调整的必要。这些尝试汇聚为下述基本准则。

\medskip
\noindent\textbf{关联准则}\quad 两个变量 $X$ 与 $Y$ 无混杂，当且仅当每一个不受 $X$ 影响的变量 $Z$ 或者
\begin{enumerate}[label=(\textup{U}\textsubscript{\arabic*}),leftmargin=3.0em]
\item 与 $X$ 无关联，或者
\item 在给定 $X$ 时与 $Y$ 无关联。
\end{enumerate}

这一准则及其若干变体与衍生物（往往回避“仅当”那一半），几乎见于每一本流行病学教科书（Schlesselman 1982; Rothman 1986; Rothman and Greenland 1998），也几乎见于每一篇讨论混杂的文章。事实上，该准则在文献中已根深蒂固，以致作者们（如 Gail 1986; Hauck et al. 1991; Becher 1992; Steyer et al. 1996）常把它当作无混杂的定义，忘记了归根结底混杂之所以有用，只在于它能告诉我们效应偏倚的情况。\footnote{Hauck 等 (1991) 把基于效应的混杂定义斥为“哲学性的”，并把两种关联度量之差视为“偏倚”。Grayson (1987) 甚至声称，作为关联准则之衍生物的“参数变化法”是混杂唯一的基本定义（对 Grayson 立场的批评见 Greenland et al. 1989）。}

本节与下一节的目的是凸显关联准则及其衍生物的若干基本局限。我们将表明，关联准则既不保证效应估计无偏，也不是从无偏性要求推出的。在用例子演示混杂的统计概念与因果概念之间并无逻辑联系之后，我们将定义一个更强的无偏性概念，称为“稳定”无偏性；相对于它，一个修改过的统计准则将被证明是充分必要的。其必要的那一半将给出一个检验稳定无偏性的实用方法，而值得注意的是，它并不要求知道问题中全部潜在的混杂因子。最后我们将论证，用统计准则取代基于效应的混杂定义这一流行做法并非完全错误，因为稳定无偏性事实上正是 (i) 研究者一直（也许应当）追求的目标，也是 (ii) 统计准则所能检验的东西。

\subsection{因果定义与关联定义}

为便于讨论，我们先将无混杂的因果定义与统计定义写成数学形式。\footnote{为简单起见，我们将讨论限于未调整的混杂；涉及辅助变量测量的推广是直接的，可由 \S 3.3 得到。我们还使用“$X$ 与 $Y$ 无混杂”这一简略说法，尽管“$X$ 对 $Y$ 的效应无混杂”更为确切。}

\begin{definition}[无混杂：因果定义]
设 $M$ 是数据生成过程的因果模型——即对每个被观测变量之取值如何被决定的形式描述。以 $P(y\mid \dox{x})$ 记在假设性干预 $X=x$ 之下响应事件 $Y=y$ 的概率，按 $M$ 计算。我们称 $X$ 与 $Y$ 在 $M$ 中无混杂，当且仅当
\begin{equation}
P(y\mid \dox{x}) = P(y\mid x) \tag{6.10}
\end{equation}
对其各自论域中的所有 $x$ 与 $y$ 成立，其中 $P(y\mid x)$ 是由 $M$ 生成的条件概率。若 \textup{(6.10)} 成立，我们说 $P(y\mid x)$ 是无偏的。
\end{definition}

就我们此处的讨论而言，我们把这个因果定义当作“无混杂”一语的含义。概率 $P(y\mid \dox{x})$ 已在第三章定义（定义 3.2.1，亦简记作 $P(y\mid \hat{x})$）；它可解释为对应于“$X$ 被随机化”的受控实验的条件概率 $P^*(Y=y\mid X=x)$。我们记得，这一概率可由因果模型 $M$ 计算：或者直接模拟干预 $\dox{X=x}$，或者（若 $P(x,s)>0$）经由调整公式（式 \textup{(3.19)}）
\[
P(y\mid \dox{x}) = \sum_{s} P(y\mid x,s)\,P(s),
\]
其中 $S$ 代表任何满足后门准则（定义 3.3.1）的变量集，无论被观测与否。等价地，$P(y\mid \dox{x})$ 可写作 $P(Y(x)=y)$，其中 $Y(x)$ 是 \textup{(3.51)} 或 Rubin (1974) 所定义的潜在结局变量。我们牢记：$\dox{\cdot}$ 算子、因而效应估计与混杂，都必须相对于特定的因果模型或数据生成模型 $M$ 来定义，因为这些概念不是统计性的，无法用联合分布来定义。

\begin{definition}[无混杂：关联准则]
设 $T$ 是问题中不受 $X$ 影响的变量集。我们称在 $T$ 在场时 $X$ 与 $Y$ 无混杂，若 $T$ 的每个成员 $Z$ 至少满足下述条件之一：
\begin{enumerate}[label=(\textup{U}\textsubscript{\arabic*}),leftmargin=3.0em]
\item $Z$ 与 $X$ 无关联（即 $P(x\mid z)=P(x)$）；
\item $Z$ 在给定 $X$ 时与 $Y$ 无关联（即 $P(y\mid z,x)=P(y\mid x)$）。
\end{enumerate}
反之，若 $T$ 中有成员 $Z$ 同时违反 (U\textsubscript{1}) 与 (U\textsubscript{2})，则称 $X$ 与 $Y$ 是混杂的。
\end{definition}

请注意，定义 6.2.2 中的关联准则并非纯统计的，因为它援引了“受……影响”这一谓词；该谓词无法从概率中识别，而依赖于因果信息。排除受处理（或暴露）影响的变量是不可避免的，且早已被认识到是任何处理效应分析（无论观测研究还是实验研究）所必需的判断性输入（Cox 1958, p.~48; Greenland and Neutra 1980）。我们将始终假定，研究者具备区分受处理 $X$ 影响的变量与不受其影响的变量所需的知识。然后我们将探究：要建立混杂的检验，还需要（若有的话）哪些额外的因果知识。

\section{关联准则如何失效}

我们将说，一个无混杂准则若是\emph{充分的}，则当它把某情形归类为无混杂时从不出错；若是\emph{必要的}，则当它把某情形归类为混杂时从不出错。定义 6.2.2 的关联准则与定义 6.2.1 的因果准则不相符的方式有若干种。关于充分性与必要性的失效，将依次讨论。

\subsection{因边缘性而失去充分性}

定义 6.2.2 的准则基于对 $T$ 中每个元素逐一检验。完全可能出现这样的情形：两个因素 $Z_1$ 与 $Z_2$ 联合地（在定义 6.2.2 的意义上）混杂 $X$ 与 $Y$，而每个因素单独却满足 (U\textsubscript{1}) 或 (U\textsubscript{2})。这之所以可能，是因为 $X$ 与 $T$ 中各成员之间的统计独立，并不保证 $X$ 与取自 $T$ 的变量组之间独立。例如，设 $Z_1$ 与 $Z_2$ 是两枚独立公平硬币的结果，各自影响 $X$ 与 $Y$。假定 $Z_1$ 与 $Z_2$ 相等时 $X$ 发生，而 $Z_1$ 与 $Z_2$ 不等时 $Y$ 发生。显然，$X$ 与 $Y$ 被数对 $T=(Z_1,Z_2)$ 严重混杂；事实上它们完全（负）相关，却彼此没有因果影响。然而 $Z_1$ 与 $Z_2$ 都不与 $X$ 或 $Y$ 关联；发现任一硬币的结果，都不会使 $X$（或 $Y$）的概率偏离其初始值 $\tfrac12$。

若试图通过把 (U\textsubscript{1}) 与 (U\textsubscript{2}) 中的 $Z$ 替换为 $T$ 的任意子集来补救定义 6.2.2，那将过于严格，因为 $X$ 与 $Y$ 的全部原因所构成的集合，若作为一个组来处理，几乎必然通不过 (U\textsubscript{1}) 与 (U\textsubscript{2}) 的检验。在 \S 6.5.2 中，我们指出若要恢复充分性，应当用哪些子集替换 (U\textsubscript{1}) 与 (U\textsubscript{2}) 中的 $Z$。

\subsection{因封闭世界假设而失去充分性}

所谓“封闭世界”假设，我指的是：我们的模型涵盖了全部相关变量，并且具体到定义 6.2.2，变量集 $T$ 由问题中全部潜在混杂因子构成。为正确归类每一个无混杂情形，关联准则要求对问题中每一个潜在混杂因子 $Z$ 都满足条件 (U\textsubscript{1}) 或 (U\textsubscript{2})。在实践中，由于研究者永远无法确定给定的潜在混杂因子集 $T$ 是否完备，关联准则会把某些混杂情形误判为无混杂。

这一局限实际上意味着：任何统计检验都注定是不充分的。由于实用检验总是只涉及 $T$ 的真子集，我们凭统计手段所能期望达到的至多是\emph{必要性}——即这样一个检验：当 (U\textsubscript{1})、(U\textsubscript{2}) 之类的准则被 $T$ 的某个任意子集违反时，它能把该情形正确地标为混杂。这一前景也未能由定义 6.2.2 实现，我们下面就来演示。

\subsection{因贫瘠代理而失去必要性}

\begin{example}
设想这样一种情形：暴露（$X$）受个人教育程度（$E$）影响，疾病（$Y$）受暴露与年龄（$A$）二者影响，而车型（$Z$）受年龄（$A$）与教育程度（$E$）二者影响。这些关系示意于图 6.3。
\end{example}

\begin{figure}[htbp]
\centering
\begin{tikzpicture}[causal,scale=0.95,every node/.style={transform shape}]
\node[vtx] (E) at (1.2,2.6) {$E$};
\node[vtx] (A) at (5.0,2.6) {$A$};
\node[vtx] (Z) at (3.1,1.5) {$Z$};
\node[vtx] (X) at (1.0,0) {$X$};
\node[vtx] (Y) at (5.2,0) {$Y$};
\draw[ed] (E)--(X);
\draw[ed] (E)--(Z);
\draw[ed] (A)--(Z);
\draw[ed] (A)--(Y);
\draw[ed] (X)--(Y);
\node[anchor=west,font=\footnotesize,align=left] at (6.1,1.7) {$X$——暴露\\
$Y$——疾病\\
$Z$——病人拥有的车型\\
$E$——教育程度\\
$A$——年龄};
\end{tikzpicture}
\caption{$X$ 与 $Y$ 无混杂，尽管 $Z$ 与二者都关联。}
\label{fig:6.3}
\end{figure}

车型变量（$Z$）违反了定义 6.2.2 的两个条件，因为：(1) 车型指示教育程度，因而与暴露变量关联；(2) 车型指示年龄，因而在暴露者与非暴露者中都与疾病关联。然而在这个例子里，$X$ 对 $Y$ 的效应并无混杂；一个人拥有的车型对其暴露或疾病都没有影响，它不过是众多无关属性之一，而这些属性经由中介与二者都关联。第三章的分析确立：确实，\textup{(6.10)} 在该模型中成立，\footnote{因为（后门）路径 $X\leftarrow E\to Z\leftarrow A\to Y$ 被 $Z$ 处的对撞箭头阻断（见定义 3.3.1）。}而且对 $Z$ 作调整一般会产生有偏结果：
\[
\sum_{z} P(Y=y\mid X=x,Z=z)\,P(Z=z) \ne P(Y=y\mid \dox{x}).
\]
于是我们看到，基于统计关联的传统准则未能识别出一个无混杂的效应，反而会诱使人们去调整错误的变量。每当我们把 (U\textsubscript{1}) 与 (U\textsubscript{2}) 施于一个\emph{贫瘠代理}（barren proxy）——即对 $X$ 或 $Y$ 没有影响、却是有此类影响之因素的代理的变量——这种失效就会发生。

读者也许认为这种失效不太严重，因为经验丰富的流行病学家很少会把一个变量当作混杂因子，除非怀疑它对 $X$ 或 $Y$ 有某种影响。尽管如此，对代理变量作调整在流行病学中是一种流行做法，应当极为谨慎（Greenland and Neutra 1980; Weinberg 1993）。为规范这种谨慎，必须修改关联准则，把贫瘠代理排除出检验集 $T$。由此得到下述修改准则，其中 $T$ 仅由（因果地）影响 $Y$（可能经由 $X$）的变量构成。

\begin{definition}[无混杂：修改的关联准则]
设 $T$ 是问题中不受 $X$ 影响、但可能影响 $Y$ 的变量集。我们称 $X$ 与 $Y$ 因 $T$ 的在场而无混杂，当且仅当 $T$ 的每个成员 $Z$ 满足定义 6.2.2 的 (U\textsubscript{1}) 或 (U\textsubscript{2})。
\end{definition}

Stone (1993) 与 Robins (1997) 提出了定义 6.2.2 的另一些修改，它们避开了贫瘠代理造成的问题，却不要求人们判断某变量是否对 $Y$ 有影响。他们不把 $T$ 限制为 $Y$ 的潜在原因，而让 $T$ 保持为不受 $X$ 影响的全部变量之集，\footnote{另一种做法是把 $T$ 限于任何足以控制混杂的变量集 $S$：$P(y\mid \dox{x})=\sum_s P(y\mid x,s)P(s)$。然而，我们同样永远无法确定模型中的被测量变量是否包含这样一个集合，也无法确定 $T$ 的哪些子集具有这一性质。}转而要求 $T$ 由两个不相交的子集 $T_1$ 与 $T_2$ 构成，使得
\begin{enumerate}[label=(\textup{U}\textsubscript{\arabic*}*),leftmargin=3.4em]
\item $T_1$ 与 $X$ 无关联；且
\item $T_2$ 在给定 $X$ 与 $T_1$ 时与 $Y$ 无关联。
\end{enumerate}

在图 6.3 的模型中，例如，取 $T_1=A$、$T_2=\{Z,E\}$ 即满足条件 (U\textsubscript{1}*) 与 (U\textsubscript{2}*)，因为（用 d-分离检验）$A$ 与 $X$ 独立，而 $\{E,Z\}$ 在给定 $\{X,A\}$ 时与 $Y$ 独立。

关联准则的这一修改进一步纠正了与边缘性有关的问题（见 \S 6.3.1），因为 (U\textsubscript{1}*) 与 (U\textsubscript{2}*) 把 $T_1$ 与 $T_2$ 当作复合变量来处理。然而，这一修改仍不足以恢复必要性。由于集合 $T=(T_1,T_2)$ 必须包含一切不受 $X$ 影响的变量（见注 13），而实用检验又限于 $T$ 的真子集，我们不能仅仅根据 (U\textsubscript{1}*)、(U\textsubscript{2}*) 的失效（如 \S 6.3.2 所述）就断定混杂存在。因此，这一准则同样不足以作为实际检测混杂的基础。

我们现在讨论另一项根本局限：我们凭统计手段检测混杂的能力所受的限制。

\subsection{因偶然抵消而失去必要性}

这里我们给出一个不含贫瘠代理的情形，其中 $X$ 对 $Y$ 的效应 (i) 在 \textup{(6.10)} 的意义上无混杂，但 (ii) 按定义 6.3.2 的修改关联准则却是混杂的。

\begin{example}
考虑由下述线性方程定义的因果模型
\begin{gather}
x = \alpha z + \varepsilon_1, \tag{6.11}\\
y = \beta x + \gamma z + \varepsilon_2, \tag{6.12}
\end{gather}
其中 $\varepsilon_1$ 与 $\varepsilon_2$ 是相关的不测量，$\mathrm{cov}(\varepsilon_1,\varepsilon_2)=r$，而 $Z$ 是一个与 $\varepsilon_1$、$\varepsilon_2$ 都不相关的外生变量。与该模型相联系的图见图 6.4。$X$ 对 $Y$ 的效应用路径系数 $\beta$ 量化，它给出 $E(Y\mid \dox{x})$ 随 $x$ 单位变化的变化率。\footnote{见 \S 3.5—3.6，或 \S 5.4.1 中的 \textup{(5.24)}。}
\end{example}

\begin{figure}[htbp]
\centering
\begin{tikzpicture}[causal,scale=0.95,every node/.style={transform shape}]
\node[vtx] (Z) at (2.9,2.6) {$Z$};
\node[vtx] (X) at (1.0,0.9) {$X$};
\node[vtx] (Y) at (4.8,0.9) {$Y$};
\draw[ed] (Z)--(X) node[midway,above left=-1pt] {$\alpha$};
\draw[ed] (Z)--(Y) node[midway,above right=-1pt] {$\gamma$};
\draw[ed] (X)--(Y) node[midway,above] {$\beta$};
\draw[bid] (X) to[bend right=32] (Y);
\node[font=\small] at (2.9,0.02) {$r$};
\end{tikzpicture}
\caption{$Z$ 与 $X$ 和 $Y$ 都关联，然而 $X$ 对 $Y$ 的效应无混杂（当 $r=-\alpha\gamma$ 时）。}
\label{fig:6.4}
\end{figure}

不难证明（假定变量已标准化），$Y$ 对 $X$ 的回归给出
\[
y = (\beta + r + \alpha\gamma)\,x + \varepsilon,
\]
其中 $\mathrm{cov}(x,\varepsilon)=0$。于是，只要等式 $r=-\alpha\gamma$ 成立，回归系数 $r_{YX}=\beta+r+\alpha\gamma$ 便是 $\beta$ 的无偏估计，这意味着 $X$ 对 $Y$ 的效应无混杂（无须作任何调整）。然而，关联条件 (U\textsubscript{1}) 与 (U\textsubscript{2}) 都被变量 $Z$ 违反：$Z$ 与 $X$ 关联（若 $\alpha\ne 0$），且在给定 $X$ 时与 $Y$ 条件关联（除 $\gamma$ 的某些特殊取值使 $\rho_{YZ\cdot X}=0$ 之外）。

这个例子表明，无偏性条件（定义 6.2.1）并不蕴含定义 6.3.2 的修改准则。关联准则可能把某些无混杂情形误判为混杂；更糟的是，对虚假混杂因子（本例中的 $Z$）作调整，会把偏倚引入效应估计。\footnote{注意，定义 6.3.2 的 Stone–Robins 修改在本例中也会失效，除非我们能测量造成 $\varepsilon_1$ 与 $\varepsilon_2$ 之间相关的那些因素。}
